% latexmk -cd -lualatex works/A_9/appendix.tex
% latexmk -c -cd works/A_9/appendix.tex

\documentclass{ees}

\newlist{lyricslist}{description}{1}
\setlist[lyricslist]{
  partopsep=0pt,
  labelindent=0em,
  labelwidth=4.5em,
  leftmargin=4.5em,
  labelsep=0pt,
  first=\ltseries,
  font=\normalfont\itshape\ltseries,
  style=multiline
}

\newenvironment{lyrics}[1]{%
  \subsection{#1}\nopagebreak%
  \begin{lyricslist}%
  \let\voice\item%
}{%
  \end{lyricslist}%
}

\begin{document}

\clearscrheadfoot

\chapter{Lyrics}

\textit{Roles:} St: Magdalena (S) · St: Michael (A) · Judas (T) · Lucifer (T) · Petrus (T) · Judex (B) · Coro (SATB)

\section{Actus primus}

\begin{lyrics}{Scena prima}
  \voice[Judex]
  Weill dan der Sünder Geil\\
  sich mehrt und überhäuffet,\\
  und die erboßte Weld\\
  veracht ihr Seelenheyll,\\
  sich auf Barmherzigkheit\\
  und meine Langmuth ſteiffet,\\
  ſo ſoll mein Zorn ergrimmen,\\
  ich will ſie als ein Hafnerwerkch\\
  zerſchmettern und zertrümmern.\\
  Wollan, bereithet euch zur Stelle,\\
  ihr Siben Engelsgeiſter!\\
  Und höret mit Erſtaunen,\\
  was ich nun anbefehle,\\
  eur Gott, eur Herr und Meiſter,\\
  man blaſe auf das erſchröckhlichſte\\
  die große Zornpoſaunen.
\end{lyrics}

\begin{lyrics}{Scena secunda}
  \voice[Chorus]
  Ach! Ô unerhörte Sach!\\
  Ach! dergleichen nie geſchehen,\\
  es muß die gantze Weld\\
  verſchröckhen, untergehen,\\
  das Feur mit Blueth vermengt\\
  fallt heufig her von oben,\\
  ach! ach! ach hätten wür doch nicht\\
  die Bueß ſo lang verſchoben!

  \voice[Judex]
  Diß iſt ein kleiner Anfang\\
  der vorbeſtimmten Straffen,\\
  ich werd ſie alle noch zugleich\\
  in meinen Zorn hinraffen.

  \voice[Chorus]
  Ô waß ein Ungeheur!\\
  Ein Berg mit puren Feur\\
  bewegt ſich ungefehr,\\
  und ſtürtzet in das Meer.

  \voice[Judex]
  Nunmehro will ich lachen\\
  zu ihren Untergang,\\
  die Erd ſoll nun erkrachen\\
  ob der Poſaunen Klang.

  \voice[Chorus]
  Ô Gott! waß würd noch werden?\\
  Eß fallt ein großer Stern\\
  herabwerths auf die Erden,\\
  der würd uns all verzehrn.

  \voice[Judex]
  Die Weld ſoll noch verſchmachten\\
  was über ſie wird kommen,\\
  dieweil ſie vor nur lachten\\
  und ſpotteten der Frommen.

  \voice[Chorus]
  Hilff Himmel waß iſt diſes?\\
  Die Sonn verlihret ihren Schein,\\
  der Mond zieht ſeine Strahlen ein,\\
  kein Menſch kan mehr beſtehen,\\
  wür müßen unfählbar\\
  noch all vor Forcht vergehen.

  \voice[Judex]
  So ſollen alle Plagen,\\
  die nur erdenckhlich ſein,\\
  auf ſie hinunter fallen\\
  und ſammentlich zermallen,\\
  ſie mögen immer klagen,\\
  ich ſpotte ihrer Peyn.

  \voice[Chorus]
  Ô Jammer! Elend! Noth!\\
  waß unerhörter Schröckhen!\\
  will dan der grechte Gott\\
  ſein ſchwäre Hand außſtreckhen?\\
  waß ungeheure Thier\\
  gleich denen Scorpionen,\\
  die ſonſt in Abgrund wohnen,\\
  begeben ſich herfür.

  \voice[Judex]
  Man laße jene Engel,\\
  ſo in den Euphrat\\
  bißher versteckhet waren\\
  nunmehro hervorgehen,\\
  da ſolle die böſe Weld\\
  zum Überfluß erfahren,\\
  wie ſie verdienet hat\\
  vor ihre Laſtermängel.

  \voice[Chorus]
  Waß grauſamvolles Kriegsheer\\
  ſtellt ſich vor uns mit Schröckhen!\\
  Ô daß uns doch die Berge\\
  und Bühel gleich bedecken.\\
  Die Mäuler dißer Roßen\\
  ſeynd voller Rauch und Schwefel,\\
  ſie ſtehen ſchon bereith\\
  zu ſtraffen unſern Frevel.

  \voice[Judex]
  Wollan, nun iſt beſchloßen\\
  der Untergang der Erden,\\
  ich ſchwör bey mir,\\
  es ſoll hinführ\\
  kein Zeit zur Bueß mehr werden.
\end{lyrics}

\begin{lyrics}{Scena tertia}
  \voice[St: Michael]
  Auf, auf, ihr Todten, eylet!\\
  Ihr ſolt vor Grichte gehen.

  \voice[Chorus]
  Ô harte Donnerworth!\\
  Wer wird alda beſtehen!

  \voice[St: Michael]
  Nunmehro brechet [?] an\\
  der groß und bittre Tag,\\
  wo aller Menſchen Werth\\
  mit der gerechten Waag\\
  durchforſchet ſollen werden.\\
  Hier muß Demostenes\\
  und Cicero verſtummen,\\
  die Schönheit Helenæ\\
  mit Schamröth ſich vermummen.\\
  Ein großer Alexander\\
  verliehret Muth und Stärckhe,\\
  wan er befraget wird\\
  umb ſeine Heldenwerckhe.\\
  Ein armer Lazarus\\
  wird höher æstimiret,\\
  als ein Tarquinius,\\
  der vormahls hoch ſtolzieret.\\
  Es hilfft kein appelliren,\\
  kein Gold noch Geld spendiren,\\
  der Arme wie der Reich\\
  gild hier dem Richter gleich.
\end{lyrics}

\begin{lyrics}{Aria prima}
  \voice[St: Michael]
  Ô Menſch betracht!\\
  und nicht veracht,\\
  waß dich der Glaube lehret,\\
  dan diſe Klag\\
  von jüngſten Tag\\
  hat ville ſchon bekheret.\\[1ex]
  Verzweifle nicht,\\
  thue deine Pflicht,\\
  nicht ſuche dich zu rächen,\\
  und hoff anbey,\\
  daß Gott auch ſey\\
  gerecht in ſein Verſprechen.
\end{lyrics}

\begin{lyrics}{Scena quarta}
  \voice[Judex]
  Auf auf ihr Himmels Geiſter!\\
  Begebet euch\\
  nun allzugleich\\
  in alle Theil der Erden,\\
  versamblet ihre Aſchenſtaub,\\
  damit ſie ſchnell\\
  nach mein Befehl\\
  zugleich beſeelet werden.

  \voice[St: Michael]
  Ihr Todte ſtehet auf\\
  und kommet zu Gericht,\\
  es leydet kein Verſchub,\\
  kein Proteſtiren nicht.
\end{lyrics}

\begin{lyrics}{Scena quinta}
  \voice[Petrus]
  So komme dann, empfang den Lohn,\\
  du mein gewenedeyter Leib!\\
  Laß dich nunmehr beſeelen,\\
  man wird dir vor dein Hertzenleyd\\
  ein Wolluſt Orth beſtellen.
\end{lyrics}

\begin{lyrics}{Aria secunda}
  \voice[Petrus]
  Ô meine Füeß! wie hönigſüeß\\
  würd euch die Müeh [doch] belohnet,\\
  und du mein Mund,\\
  der manche Stund\\
  die Weld zur Büeß ermahnet.\\[1ex]
  Ihr meine Glieder insgemein,\\
  nun werd ihr bald gezühret ſein\\
  mit jenen Kleid der Ehren,\\
  dorth iſt kein Leyd,\\
  kein Traurigkeit,\\
  die Freuden ſich ſtäts mehren.
\end{lyrics}

\begin{lyrics}{Scena sexta}
  \voice[St: Magdalena]
  Nun khere, meine Seel,\\
  getröſt in deinen Leib,\\
  der einſt berueffen war\\
  als ein verkhertes Weib,\\
  allein die Reu und Bueß\\
  hat alle Sünd verzehret,\\
  und mich wie dich\\
  anmüethiglich,\\
  gar herrlich ſchön verkläret.
\end{lyrics}

\begin{lyrics}{Aria tertia}
  \voice[St: Magdalena]
  Glückhſeelige Augen!\\
  die ihr mit der Laugen\\
  die Schwärtze der Sünden\\
  zu reinigen pflegt,\\
  anmüethige Thränen,\\
  nun werd ihr erkhennen\\
  wan ihr meinen Leibe\\
  wie Perlen bedeckht.\\[1ex]
  Wer ſolte nicht weinen\\
  wan Er kan erſcheinen\\
  ohn Makhel der Sünden\\
  am herlichen Tag,\\
  ein einzigen Zäher\\
  den ſchätzet man höher\\
  als alles, was immer\\
  die Weld nur vermag.
\end{lyrics}

\begin{lyrics}{Scena septima}
  \voice[Lucifer]
  Steh auf, du Höllenbraud!\\
  du ſolleſt nun erſcheinen\\
  vor Göttlichen Gericht,\\
  zu deiner größten Schand\\
  wirſtu mit allen Peynen\\
  nunmehro wohl gezücht.

  \voice[Judas]
  Verfluechte Schlangenbrueth,\\
  ſoll ich dich nun beziehen,\\
  die du dir haſt\\
  durch Sündenlaſt\\
  dein eygnes Unglickh ſelbſt gemacht,\\
  in Zeit und ewigs Leyd gebracht,\\
  du würſt der Höllenglueth\\
  hinführo nicht entflüehen.

  \voice[Lucifer]
  Ja ja, wie du geſprochen,\\
  ſo bleibt es auch darbey,\\
  dan deine falſche Reu\\
  wird ewiglich gerochen.

  \voice[Judas]
  Höll! Teuffel und alle Todstormenten,\\
  ach ſo helffet mir doch gleich\\
  mein Jammerleben endten!\\
  Wo iſt nun jener Baum,\\
  da ich mich vor erhanckhet?\\
  Iſt dan kein Blatz noch Raum\\
  mich irgends zu verhellen,\\
  in Sicherheit zu ſtellen?\\
  Ô daß ich doch nur wäre\\
  in Abgrund gleich verſenckhet.
\end{lyrics}

\begin{lyrics}{Aria quarta}
  \voice[Judas]
  Vermaledeyter Geitz!\\
  der du mich allerſeits\\
  mit Schmertzen überheuffet,\\
  gemacht daß ich [bereüths] verlohrn.\\[1ex]
  Ô daß ich doch zur Stund\\
  mich ſelbſt vernichten kunt!\\
  Ô hätt man mich erſeuffet\\
  gleich da ich war gebohrn.
\end{lyrics}

\begin{lyrics}{[Scena sine numero]}
  \voice[Lucifer]
  Nun iſt mein Neyd geſtillet,\\
  weil ich mein finſters Reich\\
  ſo woll beſetzet ſehe,\\
  und zahlreich angefüllet.\\
  Hier wird die Ehrſucht Jezabels\\
  mit Schmach und Schand erfüllet,\\
  der Durſt des trunckhnen Holofern\\
  mit gſchmoltznen Ertz erſtillet,\\
  der Phariseer Neid und Haß\\
  wird da in aller Ubermaß\\
  gefoltert und gequället,\\
  der Putiphar ſein geilen Weib\\
  iſt ſchon vor ihren zarten Leib\\
  ein Schwefelbad beſtellet.\\
  Vor alle Laſter insgemein,\\
  die immer zu erdenckhen,\\
  wird ein beſonder Marter ſeyn,\\
  die Leib und Seele kränkhen.
\end{lyrics}

\begin{lyrics}{Aria quinta}
  \voice[Lucifer]
  Waß achte ich der Peyn,\\
  ich pflege nur zu lachen,\\
  weil ich den größten Gnuß,\\
  dem Himmel zum Verdrus\\
  kan ich vergnüeget ſeyn\\
  in meinen Höllen Rachen.\\[1ex]
  Mein Abſicht ware nur\\
  die Seelen zu verblenden,\\
  durch mancherley Betrug\\
  mit falſcher Liſt und Lug\\
  die menſchlich Creatur\\
  mir ewig zu verpfänden.
\end{lyrics}

\begin{lyrics}{Scena octava}
  \voice[Judas]
  Ô daß doch diſer Tag\\
  zu unſern Troſt und Glickh\\
  in einen Augenblickh\\
  verwandlet möge werden!

  \voice[St: Michael]
  Nein nein, du irreſt weith\\
  ſambt deinen Mitgeferdten,\\
  hat Josue durch ſein Befehl\\
  die Sonne aufgehalten,\\
  damit er möchte in der Stell\\
  der Feunde Köpff zerſtalten.\\
  Wie mehr will ſich gezimmen,\\
  den der erzöhrnte Gott\\
  zu ſeiner größten Ehr,\\
  dem Sünder noch vill mehr\\
  zu ſeinem Hon und Spott,\\
  den Tag des Zohrn und Grimmen\\
  ein lange Zeit beyſetze,\\
  auf daß er ſich\\
  recht grauſamlich\\
  an ſeinen Feunden leze.
\end{lyrics}

\begin{lyrics}{Scena nona}
  \voice[Judex]
  Wollan die Zeit rueckht an,\\
  daß ich mich allgmach räche,\\
  daß Urtheil jeden ſpreche\\
  wie er verdienet hat.
\end{lyrics}

\begin{lyrics}{Aria sexta}
  \voice[Judex]
  Nun ſoll ſich mein Rach außgüeßen\\
  wan ich ſehe zu den [ſeh zu meinen] Füeßen\\
  meine Feund darnider falln\\
  und daß Zettergſchrey erſchalln.\\[1ex]
  Millionen Donnerkeule,\\
  die wie lauther Zohrenpfeyle,\\
  ſollen ſie zu Boden ſchlagn\\
  und in Grund der Höllen jagn.
\end{lyrics}

\begin{lyrics}{Scena decima}
  \voice[St: Michael]
  Bereithet euch nun allzugleich,
  die ihr alhier versamblet ſeyd,
  es kommet ſchon der Richter an
  in großer Krafft und Herrlichkeit.

  \voice[St: Magdalena]
  Ô wie bang iſt mein Gemüthe!

  \voice[St: Petrus]
  Mir erſtarret das Geblüethe.

  \voice[Judas]
  Wan ſich der Grechte ſorgt
  vor dem Richterſtuhl zu ſtehen,
  wie wird es dan endlich mir
  und meines gleichen gehen?

  \voice[Lucifer]
  Ich werde ſchon die Zunge ſpitzen,
  all und jede anzuklagn,
  und des Richters Zorn erhitzen,
  daß ihr ſolt vor Forcht und Angſt verzagn.
\end{lyrics}

\begin{lyrics}{Chorus}
  \voice[Chorus]
  So iſt dan diſer Tag
  des Jammers und der Plag
  nun würckhlich angeruckhet.
  Ô wie ſo manichmall
  hat uns der Sündenfall
  daß Gwißen hart gedruckhet,
  da wür zwar offt gehöret
  von diſen ſtrengen Gricht,
  doch bliben wür bethöret,
  und glaubten dannoch nicht.
\end{lyrics}

% \begin{lyrics}{}
%   \voice[]
% \end{lyrics}


\end{document}
